\chapter{Linguagens de programação}
\label{cap:linguagens}

\begin{objetivos}
Ao final deste capítulo você deve ser capaz de:
\begin{itemize}
  \item diferenciar linguagem de baixo nível de linguagem de alto nível no contexto do CLP;
  \item reconhecer as linguagens definidas pela norma e a vocação de cada uma;
  \item ler e escrever lógica de contatos: associação em série, em paralelo e mista;
  \item usar corretamente o contato normalmente aberto e o normalmente fechado;
  \item explicar por que a Lista de Instruções não deve ser a primeira linguagem ensinada hoje.
\end{itemize}
\end{objetivos}

\section{Por que a linguagem importa}

A linguagem é a ferramenta com que se gera o programa que coordena e sequencia as operações do
microprocessador. Ela determina três coisas ao mesmo tempo: quanto tempo se leva para escrever,
quem consegue ler o código depois e o quanto esse código sobrevive a uma troca de equipamento.

\index{linguagem de programação}

\section{Baixo nível e alto nível}

\begin{itemize}
  \item \textbf{baixo nível} --- linguagem de máquina (código binário ou hexadecimal) e
        \emph{assembly}, com mnemônicos como \texttt{MVI}, \texttt{OUT}, \texttt{MOV}. Cada
        processador tem seu conjunto próprio de registradores e instruções, o que torna o programa
        intransportável;
  \item \textbf{alto nível} --- linguagem próxima da comunicação humana. Exige compilador ou
        interpretador para chegar à linguagem de máquina, mas reduz muito o tempo de elaboração e
        dispensa conhecer a arquitetura do processador. O preço é o tempo de processamento maior.
\end{itemize}

A norma IEC 61131-3 padroniza linguagens de \emph{alto nível} para automação. Controladores que não
atendem à norma costumam oferecer apenas Ladder, no dialeto do próprio fabricante.

\section{As linguagens definidas pela norma}

\begin{figure}[H]
  \centering
  \includegraphics[width=0.92\linewidth]{10-classificacao-linguagens.pdf}
  \caption{Linguagens de programação de CLP: gráficas e textuais.}
  \label{fig:linguagens}
\end{figure}

\begin{table}[H]
  \centering
  \caption{Linguagens da IEC 61131-3, natureza e vocação.}
  \label{tab:linguagens}
  \begin{tabular}{p{1.4cm} p{2.6cm} p{1.6cm} p{8.0cm}}
    \toprule
    Sigla & Nome & Natureza & Para que serve melhor \\
    \midrule
    LD & \emph{Ladder diagram}, diagrama de contatos
      & Gráfica
      & Lógica discreta e intertravamento; a mais próxima do comando elétrico e a mais usada
        internacionalmente \\
    FBD & \emph{Function block diagram}, diagrama de blocos
      & Gráfica
      & Lógica combinatória e composição de blocos; permite construir blocos complexos a partir de
        blocos simples \\
    SFC & \emph{Sequential function chart}, passos e transições
      & Gráfica
      & Processos sequenciais: decompõe o problema em estados e condições de transição, mantendo a
        visão global da solução \\
    ST & \emph{Structured text}, texto estruturado
      & Textual
      & Lógica complexa: laços, condicionais, cálculo e manipulação de dados \\
    IL & \emph{Instruction list}, lista de instruções
      & Textual
      & Linguagem depreciada: ver a Seção~\ref{sec:il-depreciada} \\
    \bottomrule
  \end{tabular}
\end{table}

\paragraph{Texto estruturado (ST).} Tem base na informática: usa \texttt{IF-THEN}, \texttt{ELSE},
\texttt{FOR}, \texttt{WHILE} e por isso é recomendado para lógicas complexas. É a linguagem mais
próxima do que o aluno encontra em outras linguagens de programação.

\paragraph{Diagrama de blocos (FBD).} Representa a lógica por símbolos de porta e blocos, como na
lógica digital. É a linguagem natural quando o problema é combinar funções: comparar, calcular,
temporizar e encadear resultados.

\paragraph{Passos (SFC).} Também chamada antigamente de GRAFCET, parte do pressuposto de que o
processo tem sequencialidade. Consiste em \emph{passos} ligados a blocos de ação e
\emph{transições}: cada passo representa um estado particular do sistema, e a transição é a condição
que, quando verdadeira, desativa o passo anterior e ativa o seguinte.

\paragraph{Ladder (LD).} A linguagem de relé: contatos e bobinas, com a bobina na extremidade
direita, alimentados por duas barras verticais. Cada linha horizontal é uma sentença lógica em que
os contatos são as entradas e a bobina é a saída. É a mais simples de assimilar para quem já conhece
circuito de comando elétrico --- e é essa familiaridade que explica seu sucesso.

\subsection{Exemplos comparados}

A Figura~\ref{fig:exemplos-linguagens} coloca as linguagens lado a lado, resolvendo a mesma lógica
elementar --- $C = A \cdot \overline{B}$ --- em cada uma delas. Comparar os quadros é a forma mais
rápida de perceber a diferença de vocação: o texto estruturado e a lista de instruções descrevem a
operação passo a passo, enquanto o Ladder, o diagrama de blocos e os passos e transições mostram a
\emph{estrutura} da solução.

\begin{figure}[H]
  \centering
  \includegraphics[width=\linewidth]{10-exemplos-linguagens.jpeg}
  \caption{Exemplos comparados das linguagens da IEC 61131-3: Lista de Instruções, Texto Estruturado,
  Diagrama de Blocos de Função, Ladder e Passos e Transições, todos para a mesma lógica.}
  \label{fig:exemplos-linguagens}
\end{figure}
\fontefig{Imagem obtida na internet e reeditada com apoio de ferramenta de IA generativa (Nano Banana).}

Repare que apenas três quadros descrevem a lógica de forma gráfica, e que o exemplo em Lista de
Instruções aparece aqui como registro do que a norma já removeu --- a Seção~\ref{sec:il-depreciada}
volta ao ponto.

\section{Lendo e escrevendo Ladder}

Cada operando --- contato ou bobina --- é identificado por um endereço de memória. O endereço aparece
como nome simbólico escolhido pelo fabricante, o que é uma das razões da baixa portabilidade entre
marcas (Capítulo~\ref{cap:iec-versus-legado}).

\begin{itemize}
  \item \textbf{contato} --- conduz conforme o estado do endereço que o identifica. Os contatos
        endereçados como entrada se fecham se forem NA e se abrem se forem NF, conforme o par de
        terminais do módulo de entrada esteja acionado;
  \item \textbf{bobina} --- aciona o seu endereço: enquanto uma bobina de saída estiver atuada, o par
        de terminais correspondente no módulo de saída permanece em condução;
  \item \textbf{bit auxiliar} --- bits da memória de dados manipuláveis pelo programa, equivalentes
        aos relés auxiliares do comando elétrico. Não fazem parte da memória imagem de E/S.
\end{itemize}

\begin{figure}[H]
  \centering
  \begin{circuitikz}[scale=1]
    % Barra esquerda
    \draw (0,-0.6) -- (0,2.0);
    % Barra direita
    \draw (7.0,-0.6) -- (7.0,2.0);
    % Rung 1: contato NA em serie com contato NF, acionando a bobina
    \draw (0,1.5) to[closing switch, l=$DI1$] (2.2,1.5)
          to[opening switch, l=$\overline{DI2}$] (4.4,1.5)
          to[L, l=$DO1$] (7.0,1.5);
    % Rung 2: associacao em paralelo (logica OU) com contato em serie
    \draw (0,0.2) to[closing switch, l=$DI3$] (1.8,0.2);
    \draw (1.8,0.2) to[closing switch] (3.6,0.2);
    \draw (1.8,0.2) -- (1.8,-0.7) to[closing switch] (3.6,-0.7) -- (3.6,0.2);
    \draw (3.6,0.2) to[L, l=$DO2$] (7.0,0.2);
  \end{circuitikz}
  \caption{Lógica em Ladder: na primeira linha, associação em série (E); na segunda, associação em
  paralelo (OU).}
  \label{fig:ladder-logica}
\end{figure}

\begin{table}[H]
  \centering
  \caption{Qual contato usar na lógica, em função da chave de campo.}
  \label{tab:contato-campo}
  \begin{tabular}{l l l}
    \toprule
    Objetivo com a chave externa & Se a chave externa for & O contato no Ladder deve ser \\
    \midrule
    Ligar   & NA & NA \\
    Ligar   & NF & NF \\
    Desligar & NA & NF \\
    Desligar & NF & NA \\
    \bottomrule
  \end{tabular}
\end{table}

Pode-se usar chave externa de qualquer tipo, desde que o contato no programa seja o conveniente.
Mesmo assim, \textbf{não se deve usar chave NF para ligar nem NA para desligar}: a prática correta é
buscar a condição de \emph{falha segura}, ou seja, que a perda de sinal leve o processo ao estado
mais seguro.

O contato NF é a negação: campo fechado resulta em ausência de condução; campo aberto resulta em
condução. Essa inversão é a origem da maior parte dos erros de leitura de programa por quem está
começando.

\subsection{Associação em série, em paralelo e mista}

\begin{itemize}
  \item \textbf{série (E lógico)} --- a saída é acionada quando \emph{todos} os contatos conduzem;
  \item \textbf{paralelo (OU lógico)} --- a saída é acionada se \emph{pelo menos um} ramo conduz;
  \item \textbf{mista} --- combinações das duas, criando condições compostas.
\end{itemize}

Os três casos, escritos em Ladder e em texto estruturado:

\begin{lstlisting}[language=, caption={Lógica E: série de DI1, $\overline{DI2}$ e DI3.}, label={lst:logica-e}]
IF DI1 AND NOT DI2 AND DI3 THEN
    DO1 := TRUE;
ELSE
    DO1 := FALSE;
END_IF;
\end{lstlisting}

\begin{lstlisting}[language=, caption={Lógica OU: ramos paralelos de DI1, $\overline{DI2}$ e DI3.}, label={lst:logica-ou}]
IF DI1 OR NOT DI2 OR DI3 THEN
    DO1 := TRUE;
ELSE
    DO1 := FALSE;
END_IF;
\end{lstlisting}

\begin{lstlisting}[language=, caption={Lógica mista: DI3 em série com o paralelo de DI1 e
$\overline{DI2}$.}, label={lst:logica-mista}]
IF (DI1 OR NOT DI2) AND DI3 THEN
    DO1 := TRUE;
ELSE
    DO1 := FALSE;
END_IF;
\end{lstlisting}

A familiaridade dos exemplos não é coincidência: Ladder e texto estruturado descrevem a mesma
álgebra booleana. A diferença é de forma --- e é por isso que a norma admite as duas no mesmo
projeto.

\begin{pratica}
Usando o simulador do ambiente de programação disponível:
(a) reproduza a Figura~\ref{fig:ladder-logica} e monte a tabela-verdade das três entradas;
(b) escreva a mesma lógica em ST e compare com as listagens~\ref{lst:logica-e}
a~\ref{lst:logica-mista};
(c) inverta um dos contatos e explique, por escrito, em que condição de campo a saída passa a
comportar-se de forma insegura.
\end{pratica}

\section{Lista de Instruções: por que ela saiu de cena}
\label{sec:il-depreciada}

A Lista de Instruções (IL) é textual e sequencial, semelhante ao \emph{assembly}: uma lista de
mnemônicos executados em ordem. O material didático antigo costumava apresentá-la como linguagem
principal, com o argumento de que é ``diretamente entendida pelo controlador''. Esse argumento não
se sustenta hoje, por três razões:

\begin{enumerate}
  \item o programa em IL tem leitura difícil --- não se enxerga a lógica, apenas a sequência de
        operações;
  \item a norma \textbf{depreciou a Lista de Instruções na sua terceira edição (2013)} e a
        \textbf{removeu na quarta edição (2025)}, cuja publicação é de maio de 2025. Ensinar IL como
        linguagem de partida é ensinar a escrever código que o próprio normativo abandonou;
  \item o que ela oferecia --- acesso ao nível da máquina --- hoje é resolvido pela compilação a
        partir de ST, sem que o programador precise descer a esse nível.
\end{enumerate}

\textbf{Consequência prática para esta apostila:} IL aparece aqui como registro histórico e para
leitura de código legado, nunca como linguagem recomendada para programa novo. Quando o objetivo for
portabilidade, escreva em ST ou LD.

\begin{atencao}
A citação de edições é obrigatória quando se fala de conformidade. ``Conforme IEC 61131-3'' sem dizer
a edição é afirmação vazia: um controlador que implementa a segunda edição (2003) pode ser conforme
e, ao mesmo tempo, oferecer uma linguagem que a quarta edição já removeu.
\end{atencao}

\index{IEC 61131-3!edições}\index{Lista de Instruções}

\section{Como escolher a linguagem}

\begin{table}[H]
  \centering
  \caption{Critérios de análise das linguagens, para escolha em projeto.}
  \label{tab:criterios-linguagem}
  \begin{tabular}{p{3.2cm} p{11.4cm}}
    \toprule
    Critério & O que observar \\
    \midrule
    Forma de programação
      & \emph{Linear}: programa em um único bloco. \emph{Estruturada}: permite organizar, criar
        bibliotecas de rotinas utilitárias e dividir o programa por critério funcional, operacional
        ou geográfico \\
    Forma de representação
      & Diagrama de contatos, diagrama de blocos ou lista de instruções; a representação escolhida
        define quem consegue revisar o código \\
    Documentação
      & A documentação é recurso do ambiente de programação mais do que da linguagem, mas quanto
        mais rico em comentários e mais estruturado o código, melhor. Com muitos profissionais
        envolvidos entre concepção e manutenção, isso deixa de ser detalhe \\
    Conjunto de instruções
      & Conjunto de funções disponível: temporização, contagem, aritmética, PID, manipulação de
        dados e blocos de comunicação \\
    Portabilidade
      & Se a linguagem e os blocos usados são da norma ou do fabricante (Capítulo~\ref{cap:iec-versus-legado}) \\
    \bottomrule
  \end{tabular}
\end{table}

Programação estruturada é o que permite divisão por área, biblioteca de rotinas e manutenção por
outra pessoa que não o autor. Em curso técnico, esse é o critério que mais pesa no médio prazo: o
programa será mantido por quem não o escreveu.

\begin{exercicios}
\begin{enumerate}[label=\alph*)]
  \item Explique a diferença entre linguagem de baixo e de alto nível e dê um exemplo de cada no
        contexto do CLP.
  \item Associe cada situação à linguagem mais adequada e justifique: (a) intertravamento de uma
        prensa; (b) cálculo de conversão de unidade com laço; (c) sequência de partida de uma
        máquina de envase; (d) encadeamento de blocos de cálculo.
  \item Escreva a equação booleana e o trecho em ST para uma saída acionada por \texttt{DI1} em
        série com o paralelo de \texttt{DI2} e \texttt{DI3}.
  \item Preencha a Tabela~\ref{tab:contato-campo} para uma botoeira de emergência que deve desligar o
        processo, e explique a escolha pelo critério de falha segura.
  \item Um programa usa uma chave NF de campo para \emph{ligar} um motor. Explique o risco e como
        corrigir.
  \item Explique, com suas palavras, por que o contato NF é chamado de negação e por que isso induz
        erro de leitura.
  \item Cite as duas edições da norma citadas neste capítulo e o que cada uma fez com a Lista de
        Instruções.
  \item Se um fabricante afirma que seu controlador ``é conforme à IEC 61131-3'', que pergunta você
        faz antes de aceitar a afirmação?
\end{enumerate}
\end{exercicios}
